Beyond this spatial consideration, cross-dataset generalisation was configuration-dependent: all conditions achieved higher macro-$F_1$ on Raja than on Cao2018, but the magnitude of this difference varied by configuration. Proposed-Mean and Proposed-Med showed gaps of $+0.0712$ and $+0.0758$, respectively, which were close to the $+0.0714$ gap obtained by BLINKER-concat. By contrast, MNE-annot showed the largest gap at $+0.1037$ and was therefore the least consistent across the two corpora. This pattern indicates that the lower scores on Cao2018 are more plausibly attributable to differences between the corpora than to overfitting specific to the proposed thresholding procedure. Cao2018 was collected as an open multichannel sustained-attention driving dataset~\citep{cao2019multichannel}, and differences in acquisition hardware, electrode montage, and recording conditions can alter blink amplitude and morphology across datasets. The proposed configurations nevertheless retained a cross-dataset gap comparable to that of the established blink-specific baseline~\citep{kleifges2017blinker}. Their advantage was therefore not confined to Raja, although absolute performance remained dependent on the recording corpus.
